\section{Witness-preserving NP completeness}
\label{sec:cook}

The classical starting point is Cook's theorem on polynomially bounded
nondeterministic computations and propositional problems \cite{Cook1971}.
Cook's original paper uses polynomial oracle reducibility and tautologies.
Here we prove a modern many-one, witness-preserving circuit and CNF version,
and make its arithmetic compilation explicit.

\subsection{Circuits and normalized SAT instances}

A Boolean circuit is a finite list of gates in topological order. Input
variables precede internal gates; every wire of a gate refers to an input or
an earlier gate. The allowed internal operations are constants, NOT, AND,
and OR; OR may instead be expanded using NOT and AND. A designated wire is
the output. The circuit code includes its number of inputs and its gate
list. The validity test checks all index bounds. Evaluation fills a Boolean
array from left to right, so it is primitive recursive. Ordinary sequential
bit operations evaluate it in polynomial time in its serialization length;
a simple implementation can scan the previous wire table for every lookup.

A CNF is a finite list of clauses, and a clause a finite list of signed
variable indices. For \(3\SAT\), clauses have at most three literals.
We use normalized instances: the number \(m\) of variables is at most the
encoding length, all indices lie in \(0,\ldots,m-1\), and the input is a
well-formed binary record. A malformed instance is unsatisfiable by
definition. General binary variable names can be renamed to the distinct
names appearing in the formula, in polynomial time by pairwise comparison
and scanning. This matters because the largest numerical identifier could
be exponentially larger than the formula length. The witness is one bit per
occurring variable, not one bit per number below that largest identifier.

Write \(\Check(F,a)=1\) for a well-formed assignment satisfying all clauses.
This is a bounded conjunction of bounded disjunctions of literal bits. It has
a polynomial bit-machine implementation. Thus normalized CNF-SAT and 3SAT
have NP presentations, with unused certificate bits ignored after the
assignment's required length.

\subsection{Unrolling a verifier locally}

Fix a verifier descriptor \(v\), input \(x\) of length \(n\), witness length
\(p=p(n)\), and clock \(T=q(n)\). Convert its fixed machine to the
single-tape model if necessary, enlarging its polynomial clock. Put
\(x\#y\) on the initial tape and use a window of
\[
 W=n+p+2T+4
\]
cells. The window contains the whole input and at least \(T+1\) blank cells
of margin on either side, with harmless surplus cells. No head can reach
the external boundary during the relevant updates.

A cell label is either an ordinary tape symbol or a pair consisting of a
state and a symbol, meaning that the head is at that cell. Thus the finite
label set is
\[
 \Delta=\Gamma\ \sqcup\ (Q\times\Gamma).
\]
For valid configurations, the label of cell \(j\) after one move depends
only on the previous labels at \(j-1,j,j+1\). To see this, if the head is
at \(j\) it determines the new written symbol and whether it remains; if
it is at a neighbor it may move into \(j\); if neither occurs the symbol is
unchanged and the cell has no head. There is a fixed finite function
\(F_v:\Delta^3\to\Delta\) implementing these cases. Its values on
impossible triples, such as several conflicting heads, can be chosen
arbitrarily. The induction below only reaches valid rows.

Encode labels by a fixed number of bits, or by one-hot bits. Each output
bit of \(F_v\) is a fixed finite Boolean function. Its truth-table DNF is
an explicit circuit: OR together the ANDs specifying exactly the input
tuples on which that bit is 1. The number of gates is a computable constant
depending on the finite machine \(v\).

Construct \(C_{v,x}\) in row order. At row 0 hardwire the input, separator,
blank cells and initial head; use the \(p\) certificate bits as input wires.
For every time \(t<T\) and every cell, insert a fresh copy of the local
update circuit wired to the previous three cell labels. In the final row,
OR the tests for an accepting head. Absorbing terminal states mean the final
row reports acceptance by \(T\), including earlier halting.

\begin{lemma}[Uniform circuit construction]
\label{lem:unroll}
There is a primitive-recursive circuit constructor satisfying
\[
 \operatorname{Eval}(C_{v,x},y)
   =\Acc_{\#}(e,x,y,T).
\]
For each fixed verifier descriptor its gate count is bounded by
\[
 G\le a_v(T+1)W,
\]
where \(a_v\) is an explicitly computable constant. Its serialization and
its generation take polynomial time in \(n\).
\end{lemma}

\begin{proof}
The constructor is a collection of bounded loops over rows, cells, label
bits and the fixed local truth table. Its intermediate gates are finite
arrays, and every predecessor index is calculated from earlier row and gate
indices. This supplies a primitive-recursive syntax constructor, including
its validity proof.

Induct on \(t\). The invariant says that, for all \(j<W\), the Boolean
label in row \(t\) equals the corresponding machine cell label. It also
records that the head is unique and stays inside the allocated window.
The initial row is the definition of the machine input. At the successor
step the three previous labels are correct, and the fixed truth table is
precisely the local transition. Each claim is a bounded finite conjunction;
Theorem~\ref{thm:compile} replaces it by a Boolean equation. The final
acceptance OR proves the displayed identity.

There are \(TW\) local blocks, plus \(O(W)\) initialization and output
gates, giving the count. Each wire index needs
\(O(\log(G+p+2))\) bits. A generator which repeatedly scans its counters
and previous output takes, for example, a fixed multiple of
\((G+p+n+1)^3\) bit operations. This intentionally loose bound follows
by counting the bounded output and index-manipulation loops. Since \(p,T\)
are fixed polynomials in \(n\), both output length and runtime are
polynomial. The polynomial constants may depend on the verifier, as
many-one hardness allows. A primitive-recursive algorithm for building the
local truth tables is uniform in the verifier code, without asserting that
this preprocessing is polynomial in that variable code.
\end{proof}

\subsection{Tseitin constraints and both witness directions}

For each internal gate introduce a fresh variable \(z\). Retain the input
variables. Use the following clauses:
\begin{center}
\begin{tabular}{@{}ll@{}}
\toprule
Gate equation & Clauses\\
\midrule
\(z=\neg a\) & \((z\lor a),\ (\neg z\lor\neg a)\)\\
\(z=a\land b\) & \((\neg z\lor a),\ (\neg z\lor b),\
 (z\lor\neg a\lor\neg b)\)\\
\(z=a\lor b\) & \((z\lor\neg a),\ (z\lor\neg b),\
 (\neg z\lor a\lor b)\)\\
\(z=1\), \(z=0\) & \((z)\), respectively \((\neg z)\)\\
\bottomrule
\end{tabular}
\end{center}
Append a unit clause asserting the output variable. Also append the
tautology \((y_i\lor\neg y_i)\) for each of the \(p\) original input
variables. These tautologies retain even unused witness bits under
occurrence-based variable normalization, without constraining their values.
There are at most \(3G+p+1\) clauses. All have width at most three. If exactly three literals
are required, repeat a literal in shorter clauses. Output length is
\(O((G+p+1)\log(G+p+2))\), and the gate-to-clause conversion is polynomial
time. A direct Boolean case analysis, at most eight cases, proves that each
gadget is equivalent to its gate equation, even if incoming wires coincide.

Let \(F_v(x)\) be the resulting CNF. Define
\(\operatorname{Ext}_v(x,y)\) to be the input bits followed by all evaluated
gate bits. Define \(\operatorname{Restr}_v(x,a)\) to be the assignment's
input-bit prefix. Both are primitive-recursive functions with polynomial
length and polynomial-time implementations for fixed \(v\).

\begin{theorem}[Witness-preserving Cook--Levin]
\label{thm:cook}
For each verifier descriptor \(v\), \ipra{} proves
\begin{align*}
 \Acc_{\#}(e,x,y,T)=1
 &\ \longrightarrow\
 \Check(F_v(x),\operatorname{Ext}_v(x,y))=1,\\
 \Check(F_v(x),a)=1
 &\ \longrightarrow\
 \Acc_{\#}(e,x,\operatorname{Restr}_v(x,a),T)=1.
\end{align*}
The construction is a polynomial many-one reduction from \(N_v\) to 3SAT.
The code constructor, clocks, witness maps, and the correctness derivations
are effective finite constructions, using only the arithmetic kernel above.
\end{theorem}

\begin{proof}
Forward, evaluated gate bits satisfy every defining gate equation and hence
every gadget. Lemma~\ref{lem:unroll} gives the output bit 1, satisfying the
final unit clause.

Backward, fix a satisfying assignment. In topological order, each gate's
clauses force its assigned value to equal the Boolean operation on its
assigned predecessors. By induction, these predecessor values already equal
the actual evaluation from the restricted input. Therefore the gate value
does too. Use as invariant the bounded conjunction over all gates already
processed. The output clause now forces the actual output to be 1, and
Lemma~\ref{lem:unroll} translates it to verifier acceptance. These arguments
do not infer a witness from a double negation; they construct and project
witnesses explicitly. The size and runtime bounds follow from the two
constructors and reduction composition.
\end{proof}

\subsection{The meaning of P versus NP inside the theory}

Choose a fixed normalized 3SAT verifier and let \(\operatorname{sat}(x)\)
be its exhaustive Boolean satisfiability function. The following are
legitimate first-order sentences:
\begin{align*}
 \mathsf{SatP}&:\ \exists d\ \forall x\ge1\ P_d(x)=\operatorname{sat}(x),\\
 \mathsf{PNP}&:\ \forall v\ \exists d\ \forall x\ge1\ P_d(x)=N_v(x).
\end{align*}
Theorem~\ref{thm:cook} and the reduction compiler prove
\(\mathsf{SatP}\leftrightarrow\mathsf{PNP}\). One direction instantiates
\(v\) with SAT. For the other, given a SAT descriptor \(d\), construct
for each \(v\) the decider obtained by compiling \(F_v\) and then running
\(d\). In fact this construction gives a primitive-recursive function
\(\operatorname{Solve}(d,v)\), so the implication includes explicit
uniform machine selection. Nothing here supplies the premise
\(\mathsf{SatP}\). P versus NP remains unresolved \cite{ClayPNP}.

There are several distinct constructive lower-bound formulations. The
negation of \(\mathsf{SatP}\) is
\(\forall d\,\neg\forall x(P_d(x)=\operatorname{sat}(x))\). The stronger
statement
\[
 \forall d\ \exists x\ P_d(x)\ne\operatorname{sat}(x)
\]
actually supplies a mismatch for each descriptor. Even though equality at
each fixed \(x\) is decidable, intuitionistic logic does not generally
convert a negated unbounded universal into an existential counterexample.
The diagonal theorem below deliberately proves the stronger form for its
exponential-time language, with a named primitive-recursive witness.

\subsection{Search and decision, conditionally}

Assume a fixed polynomial SAT decider. For a satisfiable formula with
\(m\) variables, query whether its restriction with variable 0 set to 0 is
satisfiable. Keep 0 if yes, otherwise keep 1. Continue through all variables.
At each stage the formula with the chosen prefix remains satisfiable:
the disjoint cases for the next Boolean variable exhaust its assignments,
and a negative answer to one leaves a witness in the other. After \(m\)
stages the accumulated assignment itself satisfies the original formula.
An initial satisfiability test returns failure for unsatisfiable formulas.
Use ordinary CNF-SAT here; it is also reducible to 3SAT by the already
proved verifier construction, with polynomial overhead.

The stage invariant is bounded satisfiability of a finite restriction, so
it compiles to quantifier-free induction. The procedure makes at most
\(m+1\) polynomial-length SAT queries and polynomial bookkeeping steps.
For an arbitrary polynomially balanced NP search relation, first use
Theorem~\ref{thm:cook}, find a satisfying assignment, and project it back.
Thus a SAT decider yields polynomial search algorithms with explicit
witness recovery. A converse follows by applying a total search-with-failure
solver to SAT and checking its output. These are conditional compiler
theorems, not an unproved equivalence between an arbitrary relation and an
arbitrarily supplied total function.

Complementing the pointwise reduction also proves UNSAT coNP-complete in
the corresponding presentation. To reduce formula tautology to CNF-UNSAT,
negate the formula and apply the Tseitin encoding to its evaluation circuit.
The resulting CNF is unsatisfiable exactly when the original formula is a
tautology. The reverse reduction to general formula tautology simply negates
the CNF. These finite Boolean operations
do not establish \(\NP=\coNP\).
